\subsection{The closed-boundary case of a relative Cartier cell}

\begin{proposition}[AT20: Proved; a relative closed-boundary quadratic repair]
\label{AT20}
Work over $R=\mathbb Z_{(2)}$ and $S=S_R$.
Let $X\to A,B$ be actual strict maps with torsionfree terms,
$X_0=S$, $H^0(X_n)=S$, and rational coefficient comparison
in every simplex and weight. Let $A,B$ be admissible.
Suppose actual normalized elements of weight $b\geq3$ satisfy
\[
 h\in N_1(X^{b,0}),\qquad
 x_0\in N_2(X^{b,1}),\qquad
 \partial h=\rho\in S^b,\qquad
 \partial x_0=dh .
 \tag{AT.116}
\]
Put $y_0=dx_0$ and $s=2b$.
Choose a power $N=2^e$ with
\[
 \begin{split}
 N&\geq\max\{1,\mu_A(0),\mu_B(0)\},\\
 b&\geq\max\{2N,\mu_A(1),\mu_B(1)\},\\
 s+1&\geq\max\{\lambda_A(4),\lambda_B(4)\}.
 \end{split}                                                     \tag{AT.117}
\]
Use the resolved coefficient-tail construction on
$L=I_1\langle s\rangle$, $m=2$, with the displayed resolution
generators of A-S09-22.

Then both maps extend to that resolved algebra and to its
actual A-SAT01 saturation at cutoff $s+1$.
For these choices the corrected degree-three class admits
simultaneous actual relative repairs (AT.107).
In particular its relative period is zero whenever expressed
by (AT.93).
The element $x_0$ itself need not have zero zeroth face:
the hypothesis is the specified closed boundary $dh$.
No claim that every C21 row-one boundary has this form is made.
\end{proposition}

\begin{proof}
The simplicial faces of $y_0=dx_0$ all vanish, because
$\partial x_0=dh$ and $d^2h=0$.
The faces of $\eta_0=x_0y_0$ therefore also vanish.
Thus $z_0=\gamma_2(y_0)$ and $\eta_0$ give actual linear
$K(2)$ maps, by their normalized degree-two values, and
$d\eta_0=2z_0$. These are precisely the coefficient-tail
twist data required by AT16.

We first construct one explicit common reference for the
coefficient $\rho$. Write $R_N=\Gamma_R(U)$ with
$\operatorname{wt}(U)=N$, and send $U$ to $T^{[N]}$ in $S$.
Let
\[
 G_N=S\otimes_{R_N}\operatorname{DR}(\Phi,U),
 \qquad \Phi_n=R^{n+1},
 \tag{AT.118}
\]
the standard PD simplex algebra with barycentric
zero-forms $u_0,\ldots,u_n$ of weight $N$ and
$\sum u_i=U$. There are actual maps $G_N\to A,B$:
at weight $N$ choose a normalized coefficient edge in
each target with boundary $T^{[N]}$, using (AT.117).
The normalized module of $\Phi$ is the identity disk
$R[1]\to R[0]$. That edge defines its actual linear
simplicial map to target zero-forms, carrying the marked
sum to $T^{[N]}$. The free PD differential universal
property followed by the displayed base change gives
the required strict map. Only this actual elementary
construction from C09 is used here.

Write $b=kN+r_0$, with $0\leq r_0<N$; then $k\geq2$.
Let $\rho=aT^{[b]}$ and put
\[
 c_b=\frac{(kN)!}{k!(N!)^k}\binom b{r_0},
                    \qquad \lambda=a/c_b\in R.
 \tag{AT.119}
\]
The number $c_b$ is odd. Indeed the binary digits of
$kN$ and $r_0$ do not overlap, while
$s_2(kN)=s_2(k)$ and $s_2(N)=1$ make the two-adic
valuation of the factorial quotient zero.
Define in the first two simplices of (AT.118)
\[
 h'=\lambda T^{[r_0]}\gamma_k(u_1),\qquad
 x'=-\lambda T^{[r_0]}\gamma_{k-1}(u_2)\,du_1 .
 \tag{AT.120}
\]
The endpoint values of $h'$ are $0,\rho$, by (AT.119).
Deleting vertex zero sends $du_1$ in the second expression
to $du_0=-du_1$, and sends its $u_2$ to $u_1$.
Thus $d_0x'=dh'$. Deleting vertex one kills $du_1$,
and deleting vertex two kills $\gamma_{k-1}(u_2)$,
since $k\geq2$. Consequently $x'\in N_2(G_N^{b,1})$
with the required boundary. Direct differentiation gives
\[
       dx'=\lambda T^{[r_0]}
                     \gamma_{k-2}(u_2)\,du_1du_2 .
 \tag{AT.121}
\]
In particular
\[
                    \gamma_2(dx')=0,\qquad x'\,dx'=0.
 \tag{AT.122}
\]
These equalities also follow from the vanishing of forms
above degree two in $(G_N)_2$. All these are identities
in the specified integral local algebra.

Now work in $A$, denoting the images of (AT.120) by
$h'_A,x'_A$. The difference $f_A(h)-h'_A$ is a normalized
row-zero one-cycle. Choose
\[
 q\in N_2(A^{b,0}),\qquad
                  \partial q=f_A(h)-h'_A
 \tag{AT.123}
\]
using the row-zero bound in (AT.117).
Put $x^\flat=f_A(x_0)-dq$ and $y=f_A(y_0)=dx^\flat$.
Then $\partial x^\flat=dh'_A$, and
$x^\flat-x'_A$ is a normalized row-one two-cycle.
The row-one bound supplies
\[
 \alpha\in N_3(A^{b,1}),\qquad
                  d_0\alpha=x^\flat-x'_A.
 \tag{AT.124}
\]
Set $a_3=s_0x'_A+\alpha$. Its four faces are exactly
\[
                         (x^\flat,x'_A,0,0).
 \tag{AT.125}
\]
This element is not asserted to be normalized.
However (AT.122) makes both
$\gamma_2(da_3)$ and $a_3da_3$ normalized three-simplices.
Their zeroth faces are $\gamma_2(y)$ and $x^\flat y$,
and all other faces vanish. By the normalized top/bottom
construction of (AT.110), they define actual linear
$D(3)$ maps
\[
 \widetilde z_{\rm top}=\gamma_2(da_3),\qquad
 \widetilde\eta_{\rm top}=a_3da_3,\qquad
 z^\flat=\gamma_2(y),\quad \eta^\flat=x^\flat y,
 \qquad d\widetilde\eta=2\widetilde z.
 \tag{AT.126}
\]
Once again no quadratic operation on arbitrary parameter
sums has been declared linear.

All coefficient-resolution generators have weight at least
$s+1$. The argument (AT.111)--(AT.112) therefore constructs
actual resolved $D(3)$ generator values in $A$ for the
auxiliary twist (AT.126), with the terminal equation
factoring $W^0$.
This uses only target vertical comparisons through degree
four in those weights. There is no modification of the
old algebra map $f_A:X\to A$: at this intermediate step
we are specifying auxiliary linear twist data and resolved
generator values in the target.

Using the displayed generator identities of A-S09-22,
form their normalized top expressions $B^\flat$ and
$\zeta^\flat$ with the auxiliary value $\eta^\flat$.
They satisfy $dB^\flat=2\zeta^\flat$.
Their top and bottom values give actual $D(3)$ maps,
since all resolved top generators are normalized and
the two quadratic twist expressions in (AT.126) are
normalized. At weight $2s+5$, the complex
$\operatorname{Hom}(D(3),A)$ has zero cohomology in
degrees three and two. As in (AT.115), choose its
two actual primitives and restrict to the bottom.
We obtain
\[
 dU^\flat=\zeta^\flat,\qquad
                dV^\flat=2U^\flat-B^\flat
 \tag{AT.127}
\]
as actual $K(2)$ maps in the target.

It remains to restore the original twist and the original
map on $X$. Put $\Omega=qy$. Its faces vanish because those
of $y$ do, so it is an actual $K(2)$ two-form, and
\[
                  f_A(\eta_0)=\eta^\flat+d\Omega.
 \tag{AT.128}
\]
Change each top resolved $v_0$ value by
\[
                    v_0(t)\longmapsto
                              v_0(t)-\varepsilon(t)\Omega .
 \tag{AT.129}
\]
Keep every $u_0$, every higher-resolution value, and every
bottom-kernel value unchanged.
Equation (AT.128) checks its $v_0$ differential for the
original twist. The other equations hold because
$\varepsilon\partial=0$.
These assignments, together with the unchanged original
map $f_A$ on $X$, now define an actual strict PD map
from the original resolved algebra to $A$.
They factor through its torsionfree quotient.

In the notation (AT.90), the change (AT.129) sends
$a$ to $a-T\Omega$ and $b$ to $b-T^{[2]}\Omega$.
It leaves $C,H,J,g,x,r$ unchanged. Hence the original
target expressions, after this correction, obey
\[
 B_{\rm orig}=B^\flat-2J\Omega,\qquad
 \zeta_{\rm orig}=\zeta^\flat-d(J\Omega).
 \tag{AT.130}
\]
The first identity is $T H=2J$.
For the second, use $dJ=-T^{[2]}g+Tl$ in the explicit
formula for $\zeta$ and then (AT.128); it is also
obtained by differentiating the first and cancelling
two in the torsionfree target.
Therefore
\[
                   U_A=U^\flat-J\Omega,\qquad
                   V_A=V^\flat
 \tag{AT.131}
\]
satisfy precisely the two original equations (AT.107).
This finishes the construction without any appeal to
lift-choice invariance or to a period computation at
weight $s$.

Repeat for $B$ using its own fillings (AT.123)--(AT.124).
The old source, resolved source, and final relative repair
are the same for both targets.
The coefficient $H^0$ and rational comparison of the
torsionfree resolved algebra follow from the actual
AT16 contraction, as in AT19. Its simplex-zero term is
unchanged. The bounds (AT.117) supply the target degree-zero
and degree-one comparisons above $s+1$.
A-SAT01 therefore constructs its actual saturation at that
cutoff with both extending maps. Equations (AT.131) still
give the same generator values for the relative repair
over that saturation. This proves all assertions.
\end{proof}
